\subsection{TPR at Low FPR}
\label{sec:tpr}

Figure~\ref{fig:tpr_bar} shows TPR at 0.1\% and 1\% FPR across all six domains.
These are the operationally relevant thresholds: a real privacy auditor needs near-zero
false-alarm rate while catching as many actual training samples as possible.

At \textbf{1\% FPR}, XGBoost recovers 15--40\% of members depending on domain,
compared to 3--10\% for SAMA. On Pile-CC, XGBoost identifies 37.7\% of members
at this threshold; on HackerNews it reaches 40.1\%.
The improvement over SAMA is 4--8$\times$ across domains.

At \textbf{0.1\% FPR}, XGBoost still recovers 4--18\% of members.
This threshold means one false alarm per 1,000 non-members tested.
SAMA cannot operate reliably at this resolution with $N=128$ subsets,
while the trajectory classifier does so without degradation.

\begin{figure}[t]
  \centering
  \includegraphics[width=0.98\linewidth]{Plots/fig4_tpr_bar}
  \caption{TPR at 0.1\% FPR (left) and 1\% FPR (right) by domain and method.
           XGBoost and MLP substantially outperform SAMA and all black-box methods
           at both thresholds. SAMA is absent from the 0.1\% panel because its
           score distribution does not resolve at that resolution.}
  \label{fig:tpr_bar}
\end{figure}

\subsection{AUC-ROC Overview}
\label{sec:main_auc}

Table~\ref{tab:main} and Figure~\ref{fig:auc_heatmap} give the aggregate picture.
XGBoost achieves mean AUC \textbf{0.878} ($\pm$0.046) and MLP achieves
\textbf{0.882}, both well above the SAMA grey-box baseline at 0.816 ($\pm$0.037).
The +0.062 gap over SAMA comes from capturing the full shape of the ELBO curve
rather than a scalar per-subset NLL comparison.
Black-box attacks (Loss, Zlib, Ratio) cluster between 0.625 and 0.722
and gain nothing from the diffusion structure.

\begin{figure}[t]
  \centering
  \includegraphics[width=0.98\linewidth]{Plots/fig2_auc_heatmap}
  \caption{AUC-ROC across six methods and six MIMIR domains.
           XGBoost and MLP (bottom two rows) consistently reach the highest values.
           SAMA (fourth row) sits between black-box and white-box attacks.}
  \label{fig:auc_heatmap}
\end{figure}

\begin{table}[t]
\centering
\caption{AUC-ROC and TPR@1\%FPR across six MIMIR domains.
         95\% bootstrap CIs shown for XGBoost only (1000 resamples).
         \textbf{Bold}: best AUC per domain.
         Black-box: Loss, Zlib, Ratio. Grey-box: SAMA. White-box (ours): XGBoost, MLP.}
\label{tab:main}
\small
\setlength{\tabcolsep}{4.5pt}
\begin{tabular}{l l cccccc}
\toprule
& & \multicolumn{6}{c}{\textbf{Domain (AUC-ROC)}} \\
\cmidrule(lr){3-8}
\textbf{Method} & \textbf{Type} & arXiv & GitHub & H-News & Pile-CC & PubMed & Wikipedia \\
\midrule
Loss    & black-box & 0.650 & 0.705 & 0.674 & 0.683 & 0.647 & 0.662 \\
Zlib    & black-box & 0.646 & 0.722 & 0.677 & 0.708 & 0.643 & 0.674 \\
Ratio   & black-box & 0.627 & 0.703 & 0.652 & 0.682 & 0.624 & 0.640 \\
SAMA    & grey-box  & 0.783 & 0.795 & 0.833 & 0.885 & 0.777 & 0.823 \\
\midrule
XGBoost & ours & \textbf{0.877}{\tiny[.862,.891]} & 0.801{\tiny[.782,.819]} & \textbf{0.928}{\tiny[.916,.938]} & \textbf{0.930}{\tiny[.919,.941]} & \textbf{0.842}{\tiny[.824,.859]} & 0.892{\tiny[.879,.905]} \\
MLP     & ours & 0.886 & \textbf{0.824} & 0.919 & 0.924 & 0.836 & \textbf{0.902} \\
\midrule
& & \multicolumn{6}{c}{\textbf{Domain (TPR @ 1\% FPR)}} \\
\cmidrule(lr){3-8}
& & arXiv & GitHub & H-News & Pile-CC & PubMed & Wikipedia \\
\midrule
Loss    & black-box & 3.9\% & 7.1\% & 4.2\% & 3.7\% & 3.5\% & 3.9\% \\
Zlib    & black-box & 5.3\% & 7.9\% & 6.7\% & 8.3\% & 3.8\% & 5.3\% \\
SAMA    & grey-box  & 3.4\% & 3.3\% & 4.5\% & 10.4\% & 3.3\% & 4.9\% \\
XGBoost & ours & \textbf{26.9\%} & 15.9\% & \textbf{40.1\%} & \textbf{37.7\%} & \textbf{18.2\%} & 25.7\% \\
MLP     & ours & 27.4\% & \textbf{20.1\%} & 28.6\% & 31.3\% & 16.2\% & \textbf{29.1\%} \\
\bottomrule
\end{tabular}
\end{table}


\subsection{Shadow-Model Transfer}
\label{sec:shadow_results}

\begin{table}[t]
\centering
\caption{Shadow-model transfer results by condition and domain.
         \textbf{A}: Oracle (white-box upper bound, 5-fold CV on target labels).
         \textbf{B}: Shadow-46 (full 46-dim transfer, no target labels).
         \textbf{C}: ELBO+Entropy (8-dim subset only).
         \textbf{D}: Attention-only (16-dim, negative control).
         \textbf{E}: Pruned-30 (attention dims removed).
         Mean $\pm$ std computed across six domains.}
\label{tab:shadow}
\small
\setlength{\tabcolsep}{4pt}
\begin{tabular}{l ccccc c}
\toprule
& \multicolumn{6}{c}{\textbf{AUC-ROC}} \\
\cmidrule(lr){2-7}
\textbf{Domain} & \textbf{A: Oracle} & \textbf{B: Shadow-46} & \textbf{C: ELBO+H} & \textbf{D: Attn} & \textbf{E: Pruned-30} & \textbf{SAMA} \\
\midrule
arXiv         & 0.877 & 0.858 & 0.853 & 0.494 & 0.853 & 0.783 \\
GitHub        & 0.801 & 0.755 & 0.704 & 0.542 & 0.748 & 0.795 \\
HackerNews    & 0.928 & 0.912 & 0.902 & 0.541 & 0.898 & 0.833 \\
Pile-CC       & 0.930 & 0.913 & 0.907 & 0.563 & 0.907 & 0.885 \\
PubMed        & 0.842 & 0.838 & 0.840 & 0.504 & 0.842 & 0.777 \\
Wikipedia     & 0.892 & 0.874 & 0.854 & 0.507 & 0.872 & 0.823 \\
\midrule
\textbf{Mean} & \textbf{0.878} & \textbf{0.858} & 0.843 & 0.525 & 0.853 & 0.816 \\
\textbf{Std}  & 0.050 & 0.059 & 0.074 & 0.027 & 0.057 & 0.037 \\
\midrule
& \multicolumn{6}{c}{\textbf{TPR @ 1\% FPR}} \\
\cmidrule(lr){2-7}
& \textbf{A: Oracle} & \textbf{B: Shadow-46} & \textbf{C: ELBO+H} & \textbf{D: Attn} & \textbf{E: Pruned-30} & \textbf{SAMA} \\
\midrule
arXiv         & 26.9\% &  10.4\% &  9.9\% & 1.1\% & 14.2\% &  3.4\% \\
GitHub        & 15.9\% &   7.3\% &  2.3\% & 1.4\% & 10.5\% &  3.3\% \\
HackerNews    & 40.1\% &  29.5\% & 35.3\% & 1.5\% & 28.3\% &  4.5\% \\
Pile-CC       & 37.7\% &  40.4\% & 25.7\% & 1.7\% & 32.9\% & 10.4\% \\
PubMed        & 18.2\% &   9.7\% & 22.2\% & 2.0\% & 12.0\% &  3.3\% \\
Wikipedia     & 25.7\% &  19.6\% & 19.7\% & 2.0\% & 22.3\% &  4.9\% \\
\midrule
\textbf{Mean} & \textbf{27.4\%} & \textbf{19.5\%} & 19.2\% & 1.6\% & 20.0\% & 4.9\% \\
\midrule
& \multicolumn{6}{c}{\textbf{TPR @ 0.1\% FPR}} \\
\cmidrule(lr){2-7}
& \textbf{A: Oracle} & \textbf{B: Shadow-46} & \textbf{C: ELBO+H} & \textbf{D: Attn} & \textbf{E: Pruned-30} & \textbf{SAMA} \\
\midrule
arXiv         & 12.1\% &  2.4\% &  3.2\% & 0.0\% &  5.9\% & -- \\
GitHub        &  8.3\% &  3.3\% &  0.0\% & 0.8\% &  2.9\% & -- \\
HackerNews    & 14.8\% &  5.7\% & 12.9\% & 0.4\% &  3.4\% & -- \\
Pile-CC       & 17.5\% & 23.1\% &  6.0\% & 0.6\% & 15.0\% & -- \\
PubMed        &  4.1\% &  1.3\% & 10.4\% & 0.2\% &  2.3\% & -- \\
Wikipedia     &  8.5\% &  4.1\% &  0.7\% & 0.0\% &  8.5\% & -- \\
\midrule
\textbf{Mean} & \textbf{10.9\%} & \textbf{6.6\%} & 5.5\% & 0.3\% & 6.3\% & -- \\
\bottomrule
\end{tabular}
{\footnotesize ``--'' for SAMA at 0.1\% FPR: SAMA's score distribution does not resolve below 1\% FPR reliably with $N=128$ subsets.}
\end{table}


The shadow-model attack achieves results close to the white-box oracle
across all six domains (Figure~\ref{fig:shadow}, Table~\ref{tab:shadow}).
Oracle AUC averages 0.878; Shadow-46 (Condition~B) averages 0.858, a gap of
only 0.020. More strikingly, at 1\% FPR the shadow attack recovers a mean of
19.5\% of members with no target-domain labels, compared to the oracle's 27.4\%.
The gap in absolute detection rate narrows in the domains where ELBO signal is
strongest: on Pile-CC, Shadow-46 actually exceeds the oracle's TPR@1\%FPR
(40.4\% vs.\ 37.7\%), an artefact of the cross-domain calibration.

Reducing to just 8 ELBO+entropy dimensions (Condition~C) costs only 0.015 AUC
and stays within 0.2 percentage points of Shadow-46 at 1\% FPR on average,
while extracting features roughly 6$\times$ faster.
Pruned-30 (Condition~E, no attention dims) closely matches Shadow-46 throughout.

Condition~D (attention only) collapses to AUC 0.525 and TPR@1\%FPR of 1.6\%
across domains -- essentially random. Attention patterns are domain-specific and
carry no transferable membership signal. ELBO-based signals transfer;
attention signals do not.

\begin{figure}[t]
  \centering
  \includegraphics[width=0.98\linewidth]{Plots/fig6_shadow}
  \caption{Shadow-model AUC by condition and domain. Dashed lines show the
           domain-specific SAMA baseline. Oracle (A) and Shadow-46 (B) are within
           0.020 of each other in every domain. Attn-only (D) falls to near-random.}
  \label{fig:shadow}
\end{figure}

\subsection{Domain Analysis}
\label{sec:domain}

\begin{figure}[t]
  \centering
  \includegraphics[width=0.72\linewidth]{Plots/fig3_roc_aggregate}
  \caption{Aggregate ROC curves (mean $\pm$ std across six domains).
           XGBoost and MLP dominate across the full FPR range.}
  \label{fig:roc_agg}
\end{figure}

Attack difficulty varies by domain (Figure~\ref{fig:roc_agg}).
Pile-CC (AUC 0.930, TPR@1\%FPR 37.7\%) and HackerNews (0.928, 40.1\%) are
most vulnerable: high lexical diversity means the fine-tuned model's
reconstruction advantage concentrates on its training texts.
Wikipedia and arXiv are intermediate.
GitHub (AUC 0.801, TPR@1\%FPR 15.9\%) is hardest: source code is structurally
regular, so member and non-member files produce similar ELBO values even though
the raw member gap is the largest of all domains (2.75 nats).
Per-domain ROC and PR curves are in Appendix~\ref{app:domain}.

\subsection{Feature Ablation}
\label{sec:ablation}

\begin{figure}[t]
  \centering
  \includegraphics[width=0.98\linewidth]{Plots/fig5_ablation}
  \caption{Left: mean LOSO AUC drop per feature group across six domains.
           Right: mean solo AUC per group using only that group.
           The ELBO trajectory dominates; all attention groups are near-zero.}
  \label{fig:ablation}
\end{figure}

Figure~\ref{fig:ablation} shows the leave-one-signal-out (LOSO) analysis.
Removing the ELBO trajectory group drops mean AUC by \textbf{0.130},
exceeding the entire gap between XGBoost and SAMA.
Predictive entropy is the next most useful group (LOSO drop 0.043).
All four attention groups contribute less than 0.006 each.
Used alone, the ELBO trajectory achieves solo AUC between 0.67 and 0.84
across domains, approaching SAMA's performance without any subset enumeration.

The shadow conditions reinforce this: ELBO+H (8 dims) reaches 0.843 AUC
and 19.2\% TPR@1\%FPR, while attention-only collapses to 0.525 AUC and 1.6\% TPR.
Full per-domain tables are in Appendix~\ref{app:ablation}.
